\section{RNN 的梯度问题}

\subsection{梯度消失与梯度爆炸}

RNN 面临的核心训练难题是\textbf{梯度消失}（Vanishing Gradient）和\textbf{梯度爆炸}（Exploding Gradient）。

\begin{figure}[H]
\centering
\includegraphics[width=0.95\textwidth]{slides-images/slide-050.jpg}
\caption{梯度消失与梯度爆炸问题的起源：损失 $J^{(4)}(\theta)$ 对早期隐藏状态的梯度需要经过多次矩阵相乘\protect\footnotemark}
\end{figure}
\footnotetext{Slides 第51页。}

考虑损失 $J^{(t)}$ 对早期隐藏状态 $h^{(1)}$ 的偏导数：

$$\frac{\partial J^{(t)}}{\partial h^{(1)}} = \frac{\partial h^{(2)}}{\partial h^{(1)}} \times \frac{\partial h^{(3)}}{\partial h^{(2)}} \times \cdots \times \frac{\partial h^{(t)}}{\partial h^{(t-1)}} \times \frac{\partial J^{(t)}}{\partial h^{(t)}}$$

每个 $\frac{\partial h^{(i)}}{\partial h^{(i-1)}}$ 涉及对权重矩阵 $W_h$ 和激活函数导数的乘积。当这些矩阵的\textbf{谱范数}（最大奇异值）：

\begin{itemize}
    \item \textbf{小于 1}：连续相乘后梯度\textbf{指数级衰减} $\rightarrow$ \textbf{梯度消失}
    \item \textbf{大于 1}：连续相乘后梯度\textbf{指数级增长} $\rightarrow$ \textbf{梯度爆炸}
\end{itemize}

\begin{figure}[H]
\centering
\includegraphics[width=0.95\textwidth]{slides-images/slide-055.jpg}
\caption{梯度消失的直觉：当 $\frac{\partial h^{(i+1)}}{\partial h^{(i)}}$ 的值较小时，梯度信号在反向传播过程中越来越弱\protect\footnotemark}
\end{figure}
\footnotetext{Slides 第56页。}

\subsection{梯度消失的后果}

\begin{warningbox}{梯度消失意味着无法学习长距离依赖}
当梯度消失时，来自远处时间步的梯度信号几乎为零。这意味着：
\begin{itemize}
    \item 模型无法学到``远处的词对当前预测很重要''这样的模式
    \item 隐藏状态虽然理论上包含所有历史信息，但梯度信号无法告诉模型如何更好地利用早期信息
    \item 模型的有效上下文窗口远小于序列长度
\end{itemize}
\end{warningbox}

\subsection{梯度爆炸的后果与解决方案}

\begin{figure}[H]
\centering
\includegraphics[width=0.95\textwidth]{slides-images/slide-060.jpg}
\caption{梯度爆炸导致 SGD 更新步长过大，使参数跳到损失曲面的``奇怪''位置，严重时产生 Inf 或 NaN\protect\footnotemark}
\end{figure}
\footnotetext{Slides 第61页。}

梯度爆炸的解决方案相对简单——\textbf{梯度裁剪}（Gradient Clipping）：

$$\text{if } \|\nabla_\theta J\| > \tau: \quad \nabla_\theta J \leftarrow \frac{\tau}{\|\nabla_\theta J\|} \cdot \nabla_\theta J$$

\begin{itemize}
    \item $\|\nabla_\theta J\|$：梯度的范数
    \item $\tau$：裁剪阈值（超参数）
\end{itemize}

\begin{figure}[H]
\centering
\includegraphics[width=0.95\textwidth]{slides-images/slide-062.jpg}
\caption{梯度裁剪：当梯度范数超过阈值时，等比缩小梯度，保持方向不变但限制步长\protect\footnotemark}
\end{figure}
\footnotetext{Slides 第63页。}

\begin{importantbox}{梯度裁剪的直觉}
梯度裁剪就像给 SGD 的更新步长设置了一个``速度限制''。它不改变梯度的方向（仍然朝着正确的方向走），只是在步子太大时把它缩小到安全范围内。这是一个简单但非常有效的技巧，在训练 RNN 时几乎是标配。
\end{importantbox}

\subsection{梯度消失的解决思路}

梯度消失比梯度爆炸更难解决。Manning 提到了几个方向：

\begin{enumerate}
    \item \textbf{LSTM}（Long Short-Term Memory）：通过门控机制（遗忘门、输入门、输出门）显式控制信息的保留和遗忘，是解决梯度消失最经典的方法
    \item \textbf{GRU}（Gated Recurrent Unit）：LSTM 的简化版本，用更少的参数实现类似效果
    \item \textbf{残差连接}（Residual Connections）：通过跳跃连接让梯度直接流过多层
\end{enumerate}

这些内容将在后续课程中详细讲解。

\subsection{本章小结}

\begin{itemize}
    \item 梯度消失和梯度爆炸源于反向传播中的矩阵连乘效应
    \item 梯度爆炸可以通过\textbf{梯度裁剪}有效解决
    \item 梯度消失更为棘手，导致 RNN 难以学习长距离依赖。主要解决方案包括 LSTM、GRU 等门控机制
    \item 这些问题最终推动了 Transformer 架构的诞生——通过自注意力机制直接建立任意位置之间的连接，彻底绕过了序列依赖问题
\end{itemize}

%% ========== Part 5: RNN 的其他应用 ==========

